\documentclass[letterpaper,10pt]{article}
\usepackage{cv}

% ------------------------------------------------------------
% Personal information
% ------------------------------------------------------------

\hypersetup{
  pdftitle={Severn Lortie's CV},
  pdfauthor={Severn Lortie}
}

\cvfoot{C}{\href{mailto:severn.lortie@queensu.ca}
                   {severn.lortie@queensu.ca}}

\cvfoot{L}{\href{https://www.linkedin.com/in/severn-lortie/}{linkedin.com/in/severn-lortie}}

\cvfoot{R}{October 2026}

% ------------------------------------------------------------
% Optional descriptions beneath an entry
% ------------------------------------------------------------

% Short paragraph beneath the current item.
% Usage: \cvdescription{Your description here.}
\newcommand{\cvdescription}[1]{%
  \par\nopagebreak\smallskip
  \noindent #1\par
}

% Compact bullets beneath the current item.
% Uses a separate list so it does not change the template's
% formatting for the main entries.
%
% Usage:
% \begin{cvdetails}
%   \item First contribution.
%   \item Second contribution.
% \end{cvdetails}
\newenvironment{cvdetails}
  {%
    \begin{list}{\textbullet}{%
      \setlength{\leftmargin}{1.3em}%
      \setlength{\rightmargin}{0pt}%
      \setlength{\labelwidth}{0.6em}%
      \setlength{\labelsep}{0.4em}%
      \setlength{\itemindent}{0pt}%
      \setlength{\listparindent}{0pt}%
      \setlength{\topsep}{3pt}%
      \setlength{\partopsep}{0pt}%
      \setlength{\itemsep}{2pt}%
      \setlength{\parsep}{0pt}%
    }%
  }
  {\end{list}}

\begin{document}

\name{Severn Lortie}

% ------------------------------------------------------------
% Education
% ------------------------------------------------------------

\block{education}

\begin{itemize}

\item
  \textbf{MASc in Robotics} \sep September 2025--Present\\
  Smith Engineering, Queen's University, Kingston\\\
  Supervisor: Dr. Jonathan Gammell\\
  Thesis: \textit{An Autonomous Sailboat and Informative Path Planner for Environmental Monitoring}

\item
  \textbf{Bachelor of Computer Engineering} \sep September 2020--April 2025\\
  Smith Engineering, Queen's University, Kingston
\end{itemize}

% ------------------------------------------------------------
% Research experience
% ------------------------------------------------------------

\block{research experience}

\begin{itemize}

\item
  \textbf{MASc Researcher} \sep September 2025--Present\\
  \href{https://robotic-esp.com/}{Robotic ESP Lab}, Queen's University, Kingston \\
  Supervisor: Dr. Jonathan Gammell

  \begin{cvdetails}
    \item Designed and deployed the complete control, actuation, and motion-planning stack for an autonomous sailboat
    \item Validated the system through 6 on-water field tests in Lake Ontario
    \item Achieved reliable waypoint following and optimal sail angle control, hitting 8/8 test waypoints in a summer 2026 field test
    \item Designed and evaluated a sampling-based informative path planning algorithm that integrated sorted search heuristics for better performance over existing baselines
  \end{cvdetails}

\item
  \textbf{Field Assistant Technician - Drone Search and Rescue (SAR)} \sep September 2025--Present\\
  Robotic ESP Lab, Queen's University, Kingston \\
  Supervisor: Dr. Jonathan Gammell

  \begin{cvdetails}
    \item Volunteered in support of drone SAR project at the lab. Key responsibilities include: piloting recovery boat for when drone landed, transporting and setting up equipment, and monitoring surrounding water and airspace for traffic.
    \item Obtained pleasure craft and VHF radio license to support the safety of the operations.
  \end{cvdetails}

\item
  \textbf{Undergraduate Research Assistant}
  \sep May 2025--September 2025\\
  Robotic ESP Lab, Queen's University, Kingston

  \begin{cvdetails}
    \item Programmed a software bridge to ingest data from NMEA 2000, a common protocol for marine sensors, into the Robot Operating System (ROS). This enabled the development of the autonomous sailboat for my Master's research. The code was open-sourced for other researchers to use
  \end{cvdetails}

\end{itemize}

% ------------------------------------------------------------
% Work experience
% List jobs and internships with the most recent first.
% ------------------------------------------------------------

\block{work experience}

\begin{itemize}

\item
  \textbf{Software Engineering Intern} \sep May 2023--August 2024\\
  Distributive \sep Kingston, Ontario
  
  \begin{cvdetails}
    \item Built and presented an LLM inference demo on Distributive's parallel-computing platform to Nokia's CTO and AI/HPC lead, contributing to follow-up meetings and a collaboration.  
    \item Co-developed, from scratch, the company’s Python API for its core parallel-computing platform in a two-engineer effort using an embedded JavaScript-in-Python runtime; achieved feature parity with the native JavaScript API and supported several non-trivial workloads.
    \item Designed a non-blocking WebSocket bridge in C++ from the embedded V8 workers to the Chrome DevTools Protocol; the resulting debugger was used to fix critical product bugs.
  \end{cvdetails}

% Duplicate the entry above for additional jobs or internships.

\end{itemize}

% ------------------------------------------------------------
% Engineering design teams and projects
% ------------------------------------------------------------

\block{engineering design teams and projects}

\begin{itemize}

\item
  \textbf{Ingenuity Labs Racing Team} \sep RoboRacer/F1Tenth \sep September 2026--Present\\
  Queen's University

  \begin{cvdetails}
    \item Participated in three tests of the car and worked to diagnose complex fault, intermittent failure in the motor controller
    \item Unified team workflow for track mapping and race line generation, allowing testing to be conducted without large setup overhead
    \item Worked with team to diagnose fault that caused car to crash during previous competition, enabling more reliable operation
  \end{cvdetails}

\item
  \textbf{Reducing High Performance Computing Power Consumption with a Proactive CPU Governor}
  \sep Capstone Project \sep September 2024 \sep Queen's University

  \begin{cvdetails}
    \item Worked with team of three other students to design a custom CPU frequency governor to increase the power efficiency of representative HPC workloads. Testing on a scientific proxy application called Quicksilver, the team was able to achieve a 4.6\% decrease in runtime and 7.4\% decrease in total energy consumption.
    \item My role was to design and program the governor and create the final project poster.
  \end{cvdetails}

\item
  \textbf{Queen's VEX Robotics Team}
  \sep Firmware Team \sep September 2023--August 2024 \sep  Queen's University

  \begin{cvdetails}
    \item Devised high baud rate Lidar and Rotary Encoder communication channel using RP2040
    microcontrollers for the team's robot.
    \item The team passed qualifiers and competed using the communication subsystem at the
    annual VEX U Robotics World Championship in Dallas, Texas. The subsystem ran
    without issue during trials at competition, enabling the robot to perform its localization
    and mapping tasks.
  \end{cvdetails}  

\item
  \textbf{Queen's Soft Robotics Team}
  \sep Electrical Team Lead \sep August 2022--April 2023 \sep Queen's University

  \begin{cvdetails}
    \item Led a team researching lightweight, untethered power supplies for Peano-HASEL
    actuators.
    \item Guided team members through creating circuit prototypes of high voltage step-up
    converters.
    \item Played a key role in reaching a solution for a low-cost, low-power demand converter that
    could be used for the actuators.
  \end{cvdetails}  

% Duplicate an entry above for additional teams or projects.

\end{itemize}

% ------------------------------------------------------------
% Publications
% Remove this section if you have no publications to list.
% Clearly identify accepted or submitted manuscripts.
% ------------------------------------------------------------

\block{publications}

\begin{etaremune}

\item
  Sanindie Silva, David Courtis, Jax Hodgkinson, and \textbf{Severn Lortie}.
  ``Creating a Corn Yield Estimator: A Deep Learning Approach via UAV Imagery''\\
  \textit{Canadian Undergraduate Conference on AI (CUCAI)} \sep 2023
  % Add volume, pages, and DOI as appropriate.
  % For unpublished work, explicitly state its status:
  % \sep Accepted
  % \sep Submitted

\end{etaremune}

% ------------------------------------------------------------
% Presentations and posters
% ------------------------------------------------------------

\block{presentations and posters}

\begin{itemize}

\item
  \textbf{Severn Lortie}, Meredith Scharlach, and Jonathan Gammell.
  ``The Martin 16 Autonomous Sailboat''\\
  Robotics and AI Symposium (RAIS) at Ingenuity Labs, Queen's University\\ October 2026
  \sep Poster

\item
  Sanindie Silva, David Courtis, Jax Hodgkinson, and \textbf{Severn Lortie}.
  ``Creating a Corn Yield Estimator: A Deep Learning Approach via UAV Imagery''\\
  Canadian Undergraduate Conference on AI (CUCAI) \sep May 2023
  \sep Oral presentation

\end{itemize}

% ------------------------------------------------------------
% Awards and scholarships
% ------------------------------------------------------------

\block{awards and scholarships}

\begin{itemize}

\item \textbf{Ontario Graduate Scholarship (OGS)} \sep \$15,000 \sep 2026-2027

\item \textbf{Vector Institute Scholarship in AI} \sep \$17,500 \sep 2025-2026 

\item \textbf{NSERC Undergraduate Student Research Award (USRA)} \sep \$9,000 \sep May 2025

\item \textbf{Queen's University Principal's Scholarship} \sep \$4,000 \sep 2020

\item
  \textbf{2nd Place Paper}
  \sep Canadian Undergraduate Conference in AI (CUCAI) \sep \$500 \sep May 2023

\item
  \textbf{1st Place Project}
  \sep Amazon HACKTO Hackathon \sep \$1,000 \sep September 2022
\end{itemize}

% ------------------------------------------------------------
% Teaching and mentoring experience
% ------------------------------------------------------------

\block{Teaching experience}

\begin{itemize}

\item
  \textbf{Teaching Assistant - Introduction to Computer Programming for Engineers} \sep Fall 2026\\ APSC 143 \sep Queen's University

\item
  \textbf{Teaching Assistant - Introduction to Computer Programming for Engineers} \sep Winter 2026 \\ APSC 142 \sep Queen's University

\item
  \textbf{Teaching Assistant - Fundamentals of Information Structures} \sep Fall 2024 \sep ELEC 278 \\ Queen's University

\item
  \textbf{Teaching Assistant - Data Structures and Algorithms} \sep Winter 2023 \sep MREN 178\\ Queen's University

\end{itemize}

% ------------------------------------------------------------
% Technical skills
% Keep only categories relevant to your background.
% ------------------------------------------------------------

\block{technical skills}

\begin{itemize}

\item
  \textbf{Programming:}
  C/C++, Python, ROS, Julia, MATLAB/Simulink

\item
  \textbf{Electrical design:}
  PCB design, low-level communication (CAN Bus, MODBUS, serial/RS422), LTSpice 

\item
  \textbf{Design and fabrication:}
  3D printing, laser cutting, prototyping with CAD (OnShape and Fusion 360)

\end{itemize}

% ------------------------------------------------------------
% Leadership and service
% Use this for club leadership, outreach, and service.
% Keep technical design team work in the design teams section.
% ------------------------------------------------------------

\block{leadership and service}

\begin{itemize}

\item
  \textbf{Lab Documentation Manager}
  \sep Robotic ESP Lab \sep September 2025--Present


  \begin{cvdetails}
    \item  Run the lab's Wiki and work with lab members to write high-quality documentation for their hardware and software designs
    \item Manage public dataset hosting, both by setting up and maintaining server and also by working with university IT.
  \end{cvdetails} 

\end{itemize}

\end{document}
